\unit{강연 149 — 연접인 대수적 층에 대한 쌍대성 정리들}
\sid{149}
\src{169}
\seminar{(1957년 5월)}
\paperhead{연접인 대수적 층에 대한 쌍대성 정리들}%
  {알렉상드르 그로텐디크 지음}

이하의 결과들은 세르의 \fquote{대수적 쌍대성 정리}에서 영감을 받아
1955년 겨울과 1956년 겨울에 얻은 것이다. 이 결과들은 사영공간의
코호몰로지에 관한 상당히 초등적인 결과들 [3]과,
[2]에 제시된 형태의 카르탕--에일렌베르크 호몰로지 대수를 집중적으로
사용하여 매우 간단히 증명된다.

\fsection{1}{가군층의 Ext ({[2]}, 제3장 및 제4장)}

$X$를 단위원을 갖는 환(반드시 가환일 필요는 없음)의 층 $\Osh$가 주어진
위상공간이라 하자. $\Osh$-가군층들로 이루어진 아벨 범주
$\Ccat^{\Osh}$를 생각하자. 그 대상들을 $\Osh$-가군이라고도 한다.
이 범주의 모든 대상은 단사 분해를 갖는다는 것이 알려져 있으므로, 이 범주에서
잘 알려진 형식적 성질들을 지닌 $\operatorname{Ext}$ 함자들을 정의할 수
있다. 혼동을 피하기 위해 좀 더 정확히 말하면,
$\operatorname{Hom}_{\Osh}(X;A,B)$ 또는 간단히
$\operatorname{Hom}(X;A,B)$로 $A$에서 $B$로 가는 $\Osh$-준동형들의 아벨군을
나타내고, $\Homsh_{\Osh}(A,B)$로 $A$에서 $B$로 가는 준동형들의
싹의 층을 나타내겠다($A,B\in\Ccat^{\Osh}$). 고정된
$A\in\Ccat^{\Osh}$에 대하여, 각각 아벨군들의 범주
$\Ccat$와 $X$ 위의 아벨군층들의 범주
$\Ccat^{\Zsh}$에 값을 갖는 함자 $h_A$, $\hsh_A$를 다음 식으로 정의한다:
\begin{equation}\tag{1.1}
  h_A(B)=\operatorname{Hom}_{\Osh}(X;A,B),
  \qquad
  \hsh_A(B)=\Homsh_{\Osh}(A,B).
\end{equation}

$h_A$와 $\hsh_A$는 공변 왼쪽 완전 함자이다.
그 오른쪽 유도함자를 생각하고, 각각
$\operatorname{Ext}_{\Osh}^{p}(X;A,B)$와 $\Extsh_{\Osh}^{p}(A,B)$로 나타낸다. 그러므로
정의에 따라
\begin{equation}\tag{1.2}
\left\{
\begin{aligned}
\operatorname{Ext}_{\Osh}^{p}(X;A,B)
  &= (R^{p}h_A)(B)
   = H^{p}\bigl(\operatorname{Hom}_{\Osh}(X;A,C(B))\bigr),\\
\Extsh_{\Osh}^{p}(A,B)
  &= (R^{p}\hsh_A)(B)
   = H^{p}\bigl(\Homsh_{\Osh}(A,C(B))\bigr),
\end{aligned}
\right.
\end{equation}

여기서 기호 $R^{p}$는 오른쪽 유도함자를 취하는 것을 나타내고,
$C(B)$는 $\Ccat^{\Osh}$에서 $B$의 임의의 단사 분해를 나타낸다.
다음으로 나타낸
\[
  \Gamma:\Ccat^{\Zsh}\longrightarrow\Ccat
\]
\fquote{단면} 함자를 생각하자. 그 오른쪽 유도함자는 다음으로
나타낸다는 것을 상기하자:
\[
  B\longmapsto H^{p}(X,B).
\]

\src{170}
\begin{equation}\tag{1.3}
  H^{p}(X,B)=(R^{p}\Gamma)(B)=H^{p}\bigl(\Gamma(C(B))\bigr).
\end{equation}
물론 $h_A=\Gamma\hsh_A$이다. 한편 $\hsh_A$가
단사 대상을 $\Gamma$-비순환 대상으로 보냄을 확인할 수 있다. 따라서
잘 알려진 방식으로 다음을 얻는다:

\result{명제}{1} 모든 $\Osh$-가군 $A$에 대하여 $\Ccat^{\Osh}$ 위의
코호몰로지형 스펙트럼 함자가 존재하며, 이는 등급 함자
$\bigl(\operatorname{Ext}_{\Osh}^{*}(X;A,B)\bigr)$로 수렴하고 그 초항은
\begin{equation}\tag{1.4}
  E_{2}^{p,q}(A,B)=H^{p}\bigl(X,\Extsh_{\Osh}^{q}(A,B)\bigr).
\end{equation}
이로부터 \fquote{가장자리 준동형}들과 우리가 적지 않을 5항
완전열을 얻는다.

\result{따름정리}{1} $A$가 국소적으로 $\Osh^{n}$과 동형이면, 다음 표준
동형들이 성립한다:
\begin{equation}\tag{1.5}
  \operatorname{Ext}_{\Osh}^{p}(X;A,B)
  \overset{\sim}{\longleftarrow}
  H^{p}\bigl(X,\Homsh_{\Osh}(A,B)\bigr)
\end{equation}
(스펙트럼 열의 가장자리 준동형들이 주는 동형들이다). 특히 다음 표준
동형이 성립한다:
\begin{equation}\tag{1.6}
  \operatorname{Ext}_{\Osh}^{p}(X;\Osh,B)=H^{p}(X,B).
\end{equation}

이 결과들을 사용하려면 $\Extsh_{\Osh}^{p}(A,B)$를 구체적으로
계산할 수 있어야 한다. 이 함자들은 국소적으로 계산된다.
즉 $U$가 $X$의 열린집합이면 다음이 성립한다:
\[
  \Extsh_{\Osh}^{p}(A,B)|U
  =\Extsh_{\Osh|U}^{p}(A|U,B|U),
\]
이는 단사 $\Osh$-가군의 $U$로의 제한이
단사 $(\Osh|U)$-가군이기 때문이다. 또한 고정된 $x\in X$에 대하여
다음 함자적인 준동형들이 있다:
\begin{equation}\tag{1.7}
  \Homsh_{\Osh}(A,B)_x
  \longrightarrow
  \operatorname{Hom}_{\Osh_x}(A_x,B_x),
\end{equation}
이들은 $B$에 관한 코호몰로지 함자 사이의 하나의 준동형으로 유일하게
연장된다:
\begin{equation}\tag{1.8}
  \Extsh_{\Osh}^{p}(A,B)_x
  \longrightarrow
  \operatorname{Ext}_{\Osh_x}^{p}(A_x,B_x).
\end{equation}

\result{명제}{2} $A$가 $x$의 어떤 근방에서
준동형 $\Osh^{m}\longrightarrow\Osh^{n}$의 여핵과 동형이라고 하자. 그러면 $A$가 국소적으로 유한 생성 자유 가군들로 이루어진 차수 $p+1$까지의 부분 분해를 가지면 (1.8)은 차수 $p$에서 동형이고, $A$가 의사연접이면 모든 $p$에 대하여 동형이다. 하나의 유한 표시만으로는 차수 0에서의 비교만을 얻는다.\footnote{인쇄본에는 “(1.7)”로 되어 있으나,
$p$가 나타나고 앞 문장이 있으므로 “(1.8)”이어야 한다.} 특히
$A$가 연접 $\Osh$-가군이면 그러하다 [3].

\src{171}
\result{명제}{3} $L_{*}=(L_i)$를
$\Osh$-가군 $A$의 왼쪽 분해라 하자. 여기서 각 항은 국소적으로 어떤
$\Osh^{n}$과 동형인 $\Osh$-가군이다. 그러면 $\Extsh_{\Osh}^{*}(A,B)$는
$H^{*}\bigl(\Homsh_{\Osh}(L_{*},B)\bigr)$와 동일시되고,
$\operatorname{Ext}_{\Osh}^{*}(X;A,B)$는 복합체
$\Homsh_{\Osh}(L_{*},B)$에 관한 $X$의 초코호몰로지와 동일시된다.

증명은 표준적이다. $C(B)$가 $B$의 단사 분해일 때 이중복합체
$\Homsh_{\Osh}(L_{*},C(B))$와, $\Homsh_{\Osh}(L_{*},B)$ 및
$\Homsh_{\Osh}(A,C(B))$에서 이 이중복합체로 가는 자연
준동형들을 고려한다.

마지막으로, (1.2)에서 도입한 두 $\operatorname{Ext}$ 함자는 $B$에 관한
코호몰로지 함자일 뿐만 아니라, $B$에 관해서는 공변이고 $A$에
관해서는 반변인 코호몰로지 쌍함자이기도 하다.

\fsection{2}{Ext의 합성 법칙}

이 절의 결과는 CARTIER와 YONEDA가 서로 독립적으로 얻었다. 자세한
내용은 CARTIER의 강연 [1]을 보라. $\Ccat$를 아벨 범주라 하자.
$K$와 $L$을 $\Ccat$의 두 등급 대상이라 하자.
$\operatorname{Hom}(K,L)$로 다음 등급 아벨군을 나타내기로 한다. 그 차수
$n$의 성분은 $K$에서 $L$로 가는 차수 $n$의 동차 준동형들, 즉
준동형 $K^{i}\longrightarrow L^{i+n}$의 족 $(u_i)$들로 이루어진다.

$K$와 $L$이 복합체라고 하자(구체적으로 미분 연산자의 차수는 $+1$이라
하자). $\operatorname{Hom}(K,L)$에 다음과 같이 주어지는 미분
연산자를 정의한다.
\begin{equation}\tag{2.1}
  \delta(u)=du+(-1)^{n+1}ud,\qquad\text{여기서 }n=\deg(u),
\end{equation}
이로써 $\operatorname{Hom}(K,L)$은 차수 $+1$의 미분 연산자를 갖는
복합체가 된다. 차수 $n$의 순환은 (동차 연산자로서) $d$와 반가환하는
차수 $n$의 연산자이다.\footnote{인쇄본에는 “$u$와”로 되어 있으나, 식
(2.1)에 따르면 “$d$와”이어야 한다.} 따라서
$H^{*}\bigl(\operatorname{Hom}(K,L)\bigr)$을 생각할 수 있으며, 이는
$K$와 $L$의 호모토피형들의 불변량이다. 이를 $H^{*}(K,L)$로 나타내도 좋다.

세 번째 복합체 $M$이 주어지면, 준동형의 합성은 짝지음
\[
  \operatorname{Hom}(K,L)\times\operatorname{Hom}(L,M)
  \longrightarrow\operatorname{Hom}(K,M)
\]
을 정의하며 이는 미분 연산자들과 양립한다. 따라서 코호몰로지로
넘어가면 짝지음
\begin{equation}\tag{2.2}
  H^{*}(K,L)\times H^{*}(L,M)\longrightarrow H^{*}(K,M),
\end{equation}
을 얻으며, 이를 $(u,v)\longmapsto vu$로 나타낸다. 이 짝지음들은 자명한
결합법칙을 만족한다. 특히 $H^{*}(K,K)$는 단위원을 갖는 결합적
등급환이고, $H^{*}(K,L)$ (각각 $H^{*}(L,K)$)는 이 환 위의 오른쪽
(각각 왼쪽) 등급가군이다. 차수 $0$에서 (2.2)는 복합체 준동형들의
합성으로 귀착된다. 끝으로, 복합체의 완전열
\[
  0\longrightarrow K'\longrightarrow K\longrightarrow K''\longrightarrow 0,
\]
\src{172}
에서 모든 $i$에 대하여 $K'^{i}$가 $K^{i}$의 직접 인자와 동일시되면,
$\operatorname{Hom}(K'',L)$ 등을 비롯한 군의 복합체들의 완전열이
생기며, 따라서 연결 연산자
\[
  H^{i}(K',L)\longrightarrow H^{i+1}(K'',L)
\]
를 얻는다. $L$의 완전열에 관한 연결 연산자도 같은 방식으로 정의한다.
짝지음 (2.2)는 통상적인 의미에서 이 연결 연산자들과 양립한다.

이제 $\Ccat$의 모든 대상 $A$가 단사 분해 $C(A)$를 갖는다고 하자.
이중복합체 정리의 여러 변형 중 하나를 사용하면 다음이 성립한다.
\[
  H^{*}\bigl(C(A),C(B)\bigr)
  =H^{*}\bigl(\operatorname{Hom}(C(A),C(B))\bigr)
\]
은 다음과 표준적으로 동형이다.
\[
  H^{*}\bigl(\operatorname{Hom}(A,C(B))\bigr)
  =\operatorname{Ext}^{*}(A,B).
\]
위에서 살펴본 연결 연산자들은 $\operatorname{Ext}$의 연결 준동형들을
준다. 또한 여기서 짝지음 (2.2)는 연결 준동형들과 양립하는 결합적
짝지음
\begin{equation}\tag{2.3}
  \operatorname{Ext}^{*}(A,B)\times\operatorname{Ext}^{*}(B,C)
  \longrightarrow\operatorname{Ext}^{*}(A,C).
\end{equation}
을 준다. 특히 $\operatorname{Ext}^{*}(A,A)$는 단위원을 갖는 결합적
등급환이다. (같은 방식으로 $\operatorname{Ext}$ 함자들이 임의의 함자의
유도함자들에 작용함을 보일 수 있으나, 여기서는 이 사실을 사용하지 않는다.)

고려하는 범주가 $X$ 위의
$\Osh$-가군들의 범주 $\Ccat^{\Osh}$인 경우에는 다음 짝지음을 얻는다.
\begin{equation}\tag{2.4}
  \operatorname{Ext}_{\Osh}^{p}(X;A,B)
  \times\operatorname{Ext}_{\Osh}^{q}(X;B,C)
  \longrightarrow\operatorname{Ext}_{\Osh}^{p+q}(X;A,C),
\end{equation}
이 짝지음들은 앞에서 말한 대로 구체적으로 나타낼 수 있다. 더구나 같은
전개에서 아벨군의 범주를 $X$ 위의 아벨층의
범주로, $\operatorname{Hom}$ 함자들을
$\Homsh$ 함자들로 바꾸면, 같은 형식적 성질을 가지면서
이번에는 \fquote{국소적 성격}을 띠는 다음 짝지음도 정의된다:
\begin{equation}\tag{2.5}
  \Extsh_{\Osh}^{p}(A,B)\times\Extsh_{\Osh}^{q}(B,C)
  \longrightarrow\Extsh_{\Osh}^{p+q}(A,C).
\end{equation}
이 후자의 짝지음은 (1.8)의 준동형들이
$\operatorname{Ext}$ 사이의 짝지음과 양립한다는 사실을 관찰하면 더 구체화된다.

\src{173}
끝으로 $H^{p}(X,A)$ 함자들 사이에도 곱셈
구조(컵곱)가 있음을 상기하자. 그러면 명제 1의 스펙트럼 열들이
곱셈 구조와 양립함을 알 수 있다. 더
정확히 말하면, 스펙트럼 열 $E(A,B)$와 스펙트럼 열
$E(B,C)$에서 스펙트럼 열 $E(A,C)$로 가는 짝지음이 있으며, 수렴항에서는
전역 $\operatorname{Ext}$ 사이의 짝지음으로 귀착되고,
초기항에서는 (1.4)의 둘째 항에 나타나는
컵곱과 국소 $\operatorname{Ext}$의 짝지음으로부터 유도되는 짝지음으로 귀착된다. 이에 따라
특히 다음 \fquote{에지 준동형}들은
\begin{equation}\tag{2.6}
  \operatorname{Ext}_{\Osh}^{n}(X;A,B)
  \longrightarrow H^{0}\bigl(X,\Extsh_{\Osh}^{n}(A,B)\bigr)
\end{equation}
그리고
\begin{equation}\tag{2.7}
  H^{n}\bigl(X,\Homsh_{\Osh}(A,B)\bigr)
  \longrightarrow\operatorname{Ext}_{\Osh}^{n}(X;A,B)
\end{equation}
곱셈 구조와 양립한다. 따라서
국소적으로 어떤 $\Osh^{m}$과 동형인 층들로 한정하면, 이는
컵곱을 사용하여 전역 $\operatorname{Ext}$ 사이의 합성을 완전히 구체적으로 나타내며,
여기서는 동형 (1.5)를 감안하였다.

\fsection{3}{국소 코호몰로지의 결과}

$J$를 아이디얼로 갖는 단위원이 있는 가환환 $A$를 취하자. 우리는
가변 $A$-가군 $M$에 대하여 다음 함자적 준동형들을
\begin{equation}\tag{3.1}
\left\{
\begin{aligned}
\operatorname{Ext}_{A}^{p}(A/J,M)
  &\longrightarrow
  \operatorname{Hom}_{A}\!\left(\bigwedge^{p}J/J^{2},M\otimes A/J\right),\\
\operatorname{Tor}_{p}^{A}(A/J,M)
  &\longleftarrow
  \left(\bigwedge^{p}J/J^{2}\right)
  \otimes\operatorname{Hom}_{A}(A/J,M),
\end{aligned}
\right.
\end{equation}
여기서 텐서곱과 외승은 환 $A$ 위에서 취한다. 또한
$J/J^{2}$는 사실 $A/J$-가군이고 그
외승들은 이를 $A$-가군으로 보든 $A/J$-가군으로 보든 서로 같음에 유의하자.
준동형 (3.1)의 정의는 다음 데이터의 정의로 귀착된다. 즉,
$J$의 원소들로 이루어진 임의의 계 $\underline{x}=(x_1,\ldots,x_p)$마다
다음 준동형 $\varphi_{\underline{x}}$를 정의한다:
\begin{equation}\tag{3.2}
\left\{
\begin{aligned}
\varphi_{\underline{x}}:
  \operatorname{Ext}_{A}^{p}(A/J,M)&\longrightarrow M\otimes A/J,\\
\varphi_{\underline{x}}:
  \operatorname{Hom}_{A}(A/J,M)&\longrightarrow
  \operatorname{Tor}_{p}^{A}(A/J,M),
\end{aligned}
\right.
\end{equation}
이들은 다음 조건을 만족해야 한다:

\src{174}
\begin{enumerate}
\item[\textup{i.}] $\varphi_{x_1,\ldots,x_p}$는
계 $(x_i)$에 대해 교대 $A$-다중선형이다. 여기서 $x_i\in J$이다.
\item[\textup{ii.}] $\varphi_{x_1,\ldots,x_p}$는 $x_i$ 가운데 하나라도
$J^{2}$에 속하면 $0$이다.
\end{enumerate}
실제로 ii는 i에서 따른다. $a\in J$이면 $a\varphi_{\underline{x}}=0$이기 때문이다.
이는 (3.2)에 나오는 모든 가군이
$J$에 의해 소멸된다는 사실에서 알 수 있다.

$\varphi_{\underline{x}}$를 정의하기 위해, 복합체
$K_{\underline{x}}$를 생각하자. 그 바탕 $A$-가군은 $\bigwedge A^{p}$이고,
그 미분 연산자는 $A^{p}$ 위의 선형형식으로 본
성분 $x_1,\ldots,x_p$의 $\underline{x}$에 의한 내부곱 $i_{\underline{x}}$이다.
미분 연산자의 차수는 $-1$이고, 차수들은
음이 아니며, 이 복합체의 차수 $0$ 호몰로지는
$A/(x_1A+\cdots+x_pA)$이다. $x_i$들이 $J$에 속하므로,
증대사상 $K_{\underline{x},0}\longrightarrow A/J$를 얻는다. 따라서
$K_{\underline{x}}$는 증대된 자유 복합체로 나타나며, 그 증대 가군은
$A/J$이다. 잘 알려진 방법으로 다음 준동형들이 유도된다.
\[
\operatorname{Ext}_{A}^{*}\bigl(H_0(K_{\underline{x}}),M\bigr)
\longrightarrow
H^{*}\bigl(\operatorname{Hom}_{A}(K_{\underline{x}},M)\bigr)
\]
그리고
\[
\operatorname{Tor}_{*}^{A}\bigl(H_0(K_{\underline{x}}),M\bigr)
\longleftarrow
H_{*}\bigl(K_{\underline{x}}\otimes_A M\bigr),
\]
따라서 이를 증대 준동형으로부터 유도되는 $\operatorname{Ext}$와
$\operatorname{Tor}$의 준동형들과 합성한다. 그 증대 준동형은
$H_0(K_{\underline{x}})\longrightarrow A/J$이며, 다음 준동형들을 얻는다:
\begin{equation}\tag{3.3}
\left\{
\begin{aligned}
\psi_{\underline{x}}:
\operatorname{Ext}_{A}^{*}(A/J,M)
&\longrightarrow H^{*}\bigl(\operatorname{Hom}_{A}(K_{\underline{x}},M)\bigr),\\
\psi_{\underline{x}}:
\operatorname{Tor}_{*}^{A}(A/J,M)
&\longleftarrow H_{*}\bigl(K_{\underline{x}}\otimes_A M\bigr).
\end{aligned}
\right.
\end{equation}
그런데 최고 차수 $p$에서 첫째 우변의 코호몰로지는
$M/(x_1M+\cdots+x_pM)$이고, 둘째 우변의 호몰로지는 $M$에서
모든 $x_i$에 의해 소멸되는 원소들의 집합임을 즉시 알 수 있다. $x_i$들이 $J$에 속하므로
다음 준동형들을 얻는다.
\begin{equation}\tag{3.4}
\left\{
\begin{aligned}
H^{p}\bigl(\operatorname{Hom}_{A}(K_{\underline{x}},M)\bigr)
&\longrightarrow M\otimes A/J,\\
H_{p}\bigl(K_{\underline{x}}\otimes_A M\bigr)
&\longleftarrow\operatorname{Hom}_{A}(A/J,M).
\end{aligned}
\right.
\end{equation}
(3.3)과 (3.4)의 준동형들을 합성하면, 정의하고자 했던 (3.2)의
준동형들을 얻는다. i.를 확인하는 일은 번거롭지만
어렵지는 않다.

\src{175}
\result{명제}{4} $A$를 단위원을 갖는 가환환이라 하고,
$(x_1,\ldots,x_p)$를 $A$의 원소들의 열이라 하자.
모든 $1\leq i\leq p$에 대하여, $A$를 $(x_1,\ldots,x_{i-1})$이 생성하는 아이디얼로 나눈 몫에서
$x_i$의 상이 영인자가 아니라고 하자. $J$를 $x_i$들이 생성하는 아이디얼이라 하자.
그러면 $J/J^{2}$는 $x_i$들의 표준상들을 기저로 갖는 자유 $(A/J)$-가군이고,
복합체 $K_{\underline{x}}$는 $A/J$의 자유 분해이며, 모든
$A$-가군 $M$에 대하여 차수 $p$에서 (3.1)의 준동형들은
전단사이다. 또한 $J.M=0$이면 임의의 차수 $i$에서 정의된 유사한
준동형들도 전단사이다.

(다른 모든 결론이 따르는 핵심은 $K_{\underline{x}}$의 양의 차수
호몰로지가 모두 $0$이라는 사실이며, 이는 주어진 조건 아래 잘 알려져 있다.)

\result{따름정리}{1} $A$와 $J$가 위와 같다고 하고, 더 나아가
$A$가 완전체 $k$ 위의 차원 $n$인 정칙 아핀 대수이며
$A/J$도 정칙 아핀 대수라고 하자. $\Omega^{1}(A)$와
$\Omega^{1}(A/J)$로 각각의 켈러 미분 가군을 나타내자. 그러면
국소화와 양립하는 다음 표준 동형사상이 있다.
\begin{equation}\tag{3.5}
  \operatorname{Ext}_{A}^{p}\bigl(\Omega^{n-p}(A/J),\Omega^{n}(A)\bigr)=A/J.
\end{equation}
실제로 $\Omega^{n-p}(A/J)$는 계수(階數) $1$인 사영 $(A/J)$-가군이고, 마찬가지로
$\Omega^{n}(A)$는 계수(階數) $1$인 사영 $A$-가군이다.\footnote{인쇄본에는
「계수(階數) $n$」이라고 되어 있다. 최고차 형식의 두 가군은 계수(階數)
$1$이며, 이는 아래의 $\bigwedge^{p}(J/J^{2})$와의 동일시에도 필요하다.}
따라서 좌변은 다음과 같다.
\[
  \operatorname{Ext}_{A}^{p}(A/J,A)
  \otimes\Omega^{n-p}(A/J)'\otimes\Omega^{n}(A)
\]
(여기서 기호 $'$는 $(A/J)$-가군의 쌍대를 뜻한다.) 마지막 두 인자의
텐서곱은 $\bigwedge^{p}(J/J^{2})$와 동일시되므로, 전체는
다음과 동일시된다.
\[
  \operatorname{Ext}_{A}^{p}\!\left(A/J,\bigwedge^{p}(J/J^{2})\right),
\]
따라서 명제에 의해 다음과 동일시된다.
\[
  \operatorname{Hom}_{A}\!\left(\bigwedge^{p}J/J^{2},
  \bigwedge^{p}J/J^{2}\right),
\]
즉 $A/J$와 동일시된다. 특히
$\operatorname{Ext}_{A}^{p}\bigl(\Omega^{n-p}(A/J),\Omega^{n}(A)\bigr)$에는
$A/J$의 단위원에 대응하는 특별히 정해진 원소가 있으며, 이를
\emph{$A$ 안의 아이디얼 $J$의 기본류}라 한다. (실제로 이는
훨씬 더 넓은 조건 아래에서도 정의할 수 있다.) 따름정리 1은
더 기하적이고 더 전역적인 다음 형태로 나타낼 수 있다:

\result{따름정리}{2} $X$를 대수적으로 닫힌 체 $k$ 위의 비특이
다양체라 하고, $Y$를 $X$의 닫힌 비특이 부분다양체라 하자.
$\Osh_X$를 $X$의 국소환들의 층, $\Osh_Y$를 $Y$의 국소환들의 층이라
하고, 후자를 $\Osh_X$의 몫층으로 간주하자. $X$의 차원을 $n$,
$Y$의 차원을 $n-p$라 하자.

\src{176}
$\Omega_X$(각각 $\Omega_Y$)를 $X$(각각 $Y$) 위의 정칙
미분형식의 싹의 층이라 하자. 그러면 다음 표준 동형사상들이
있다:
\begin{equation}\tag{3.6}
  \Extsh_{\Osh_X}^{p}\bigl(\Omega_Y^{n-p},\Omega_X^{n}\bigr)=\Osh_Y,
\end{equation}
또는 같은 말로,
\begin{equation}\tag{3.6 bis}
  \Extsh_{\Osh_X}^{p}\bigl(\Osh_Y,\Omega_X^{n}\bigr)=\Omega_Y^{n-p}.
\end{equation}
식 (3.6 bis)는 $Y$가 특이 다양체일 때 $\Omega_Y^{n-p}$의 정의로
삼을 수 있다. 더 정확히 말하면 다음과 같다:

\result{명제}{5} $Y$를 차원이 $q=n-p$인, 차원 $n$의 비특이
대수다양체 $X$ 안의 대수적 부분집합이라 하자. $F$를 지지집합이
$Y$에 포함되는 $X$ 위의 연접인 대수적 층으로, $L$을 $X$ 위의
국소 자유인 대수적 층으로 하자. 그러면 층
$\Extsh_{\Osh_X}^{i}(F,L)$은 $i<p$일 때 영이고, $i=p$일 때에는 다음
표준 동형사상이 있다.
\begin{equation}\tag{3.7}
  \Extsh_{\Osh_X}^{p}(F,L)
  =\Homsh_{\Osh_X}\!\left(F,
  \Extsh_{\Osh_X}^{p}(\Osh_X/J,L)\right),
\end{equation}
여기서 $J$는 $F$를 소멸시키고 영점집합이 $Y$인 $X$ 위의 임의의
아이디얼층이다. 특히 $F$가 $Y$ 위의 연접인 대수적 층이면 다음이 성립한다.
\begin{equation}\tag{3.7 bis}
  \Extsh_{\Osh_X}^{p}(F,L)
  =\Homsh_{\Osh_Y}\!\left(F,\Extsh_{\Osh_X}^{p}(\Osh_Y,L)\right).
\end{equation}
끝으로 여전히 $F$가 $Y$ 위의 연접인 대수적 층이라고 할 때, 층
\[
  E^{i}(F)=\Extsh_{\Osh_X}^{p+i}(F,\Omega_X^{n})
\]
들은 대수적 공간 $Y$를 비특이 다양체 $X$ 안에 몰입시키는 방식에
의존하지 않는다.

문제는 국소적이므로 $X$가 아핀이고 $L=\Osh_X$라고 가정할 수 있다.
이로써 문제는 가환대수의 문제, 더 나아가 국소대수의 문제로 귀착된다. 즉
$A$가 정칙 국소환이고 $M$이 지지집합의 차원이
$\leq q=n-p$인 $A$-가군일 때, $\operatorname{Ext}_{A}^{i}(M,A)=0$이 $i<p$에 대해 성립하고 다음이 성립함을 보이면 된다.
\[
  \operatorname{Ext}_{A}^{p}(M,A)
  =\operatorname{Hom}_{A}\!\left(M,\operatorname{Ext}_{A}^{p}(A/J,A)\right),
\]
여기서 $J$는 $M$을 소멸시키고 \fquote{차원}이 $\leq q$인 임의의 아이디얼이다.
첫 번째 주장에는 $q$에 대한 귀납법을 쓴다. 간단한 단계적 분해 논법으로
$M$이 $A/J$ 꼴인 경우로 귀착된다. 이어서 $J$를 더 작은 아이디얼로
바꾸고 귀납 가정과 $\operatorname{Ext}$의 완전열을 이용하면, $J$가
명제 4에서와 같은 \fquote{매개변수계}로 생성되는 경우로 귀착되며, 이때 결과는 즉시 성립한다. 앞의 결과에 따르면 $J$가

\src{177}
\fquote{차원}이 $\leq q$인 고정된 아이디얼이면, $A/J$-가군의 범주에서
정의되는 반변 함자 $E(M)=\operatorname{Ext}_{A}^{p}(M,A)$는 왼쪽
완전하다. 또한 이 함자는 직합을 직접곱으로 보내므로 다음이 쉽게 따른다.
\[
  E(M)=\operatorname{Hom}_{A}\bigl(M,E(A/J)\bigr).
\]
\footnote{인쇄본에는 $E(A)$로 되어 있다. 여기서 다루는 범주에서 이를
표현하는 대상은 $A/J$이고, (3.7)에 필요한 등식은 $E(A/J)$이다.} 끝으로
명제 5의 마지막 주장은 더 미묘하며, 여기서는 제시할 수 없는 국소
쌍대성 정리를 이용하여 $E^{i}(F)$를 내재적으로 특징짓는 데서 따른다.

\result{따름정리}{} $\omega_Y^{q}$로 층
$\Extsh_{\Osh_X}^{p}(\Osh_Y,\Omega_X^{n})$을 나타내자. 그러면
$F$가 $Y$ 위의 연접 대수적 층이면 다음 함자적 동형사상이 성립한다:
\begin{equation}\tag{3.8}
  \Extsh_{\Osh_X}^{p}(F,\Omega_X^{n})
  =\Homsh_{\Osh_X}(F,\omega_Y^{q}).
\end{equation}

\fsection{4}{부분다양체에 대응하는 코호몰로지류}

이하 전체에서 $X$는 차원 $n$이고 체 $k$ 위에 정의된 대수적 집합을
나타내며, 간단히 하기 위해 이 체는 대수적으로 닫혔다고 가정한다. 6절을
제외하면 $X$는 비특이라고 가정한다. $\Osh_X$로 $X$ 위의 국소환들의 층을,
\[
  \Omega_X^{*}=\bigoplus_p\Omega_X^{p}
\]
로 $X$ 위의 미분형식들의 싹의 층을 나타낸다.\footnote{인쇄본은
합집합 기호를 사용한다. 그러나 등급 구조는 층 $\Omega_X^p$들의 직합
이다.} $Y$가 $X$의 닫힌 부분집합이면, 그 위의 연접 대수적 층을
$Y$ 밖에서 영인 $X$ 위의 연접층과 동일시한다. 특히 $\Osh_Y$와
$\Omega_Y$도 그러하다.

\result{보조정리}{1} $F$를 $X$ 위의 연접 대수적 층으로서 그 지지집합의
차원이 $\leq n-p$인 것으로 하고, $L$을 $X$ 위의 국소 자유 연접 대수적 층으로
하자. 그러면 $\operatorname{Ext}_{\Osh_X}^{i}(X;F,L)$은 $i<p$일 때 영이고,
다음 표준 동형사상이 성립한다.
\begin{equation}\tag{4.1}
  \operatorname{Ext}_{\Osh_X}^{p}(X;F,L)
  =H^{0}\bigl(X,\Extsh_{\Osh_X}^{p}(F,L)\bigr).
\end{equation}
또한 $F$가 닫힌 부분집합 $Y$ 위의 연접 대수적 층이고, 이 닫힌집합은 $X$ 안에서
차원이 $\leq n-p$이면 다음 표준 동형사상이 성립한다.
\begin{equation}\tag{4.1 bis}
  \Extsh_{\Osh_X}^{p}(F,L)
  =\Homsh_{\Osh_X}\!\left(
    F\otimes L'\otimes\Omega_X^{n},\omega_Y^{n-p}
  \right),
\end{equation}
여기서 $\omega_Y^{n-p}$는 명제 5의 따름정리에서 정의한 $Y$ 위의 층이다
($\Omega_Y^{n-p}$와 동일시되는데, 이는 $Y$가 비특이일 때이다).

공식 (4.1)은 명제 1의 스펙트럼 열과 명제 5에서 곧바로 따른다. 공식
(3.8)에 따라 다음과 같이 쓸 수 있다.

\src{178}
\begin{align*}
\Extsh_{\Osh_X}^{p}(F,L)
&=L\otimes(\Omega_X^{n})'\otimes
  \Extsh_{\Osh_X}^{p}(F,\Omega_X^{n})\\
&=L\otimes(\Omega_X^{n})'\otimes
  \Homsh_{\Osh_X}(F,\omega_Y^{q})\\
&=\Homsh_{\Osh_X}\!\left(
  F\otimes L'\otimes\Omega_X^{n},\omega_Y^{q}\right),
\end{align*}
여기서 $q=n-p$이고, 이로부터 공식 (4.1 bis)를 얻는다.

특히 $F=\Osh_Y$, $L=\Omega_X^{p}$로 두자.
$\Omega_X^{n}\otimes(\Omega_X^{p})'=\Omega_X^{n-p}$임을 고려하면 다음 표준
동형사상을 얻는다.
\begin{equation}\tag{4.2}
  \operatorname{Ext}_{\Osh_X}^{p}(X;\Osh_Y,\Omega_X^{p})
  =\operatorname{Hom}_{\Osh_X}
  (X;\Omega_X^{n-p},\omega_Y^{n-p}).
\end{equation}
이제 간단히 하기 위해 $Y$가 비특이라고 가정하자. 그러면
$\omega_Y^{n-p}=\Omega_Y^{n-p}$이다. $\Omega_X^{n-p}$에서
$\Omega_Y^{n-p}$로 가는 자연스러운 준동형이 있으므로, 층
$\Extsh_{\Osh_X}^{p}(\Osh_Y,\Omega_X^{p})$의 표준 단면 $s_Y$를 얻는다.
$Y$의 모든 성분의 차원이 $n-p$라면 이 단면을 층
$\Extsh_{\Osh_X}^{p}(\Osh_Y,\Omega_X^{p})$의 \emph{기본 단면}이라 부르겠다.
이 단면은 (4.1)에 따라
$\operatorname{Ext}_{\Osh_X}^{p}(X;\Osh_Y,\Omega_X^{p})$의 한 원소를 정한다. 한편
자연스러운 준동형 $\Osh_X\longrightarrow\Osh_Y$는 다음 준동형을 정한다:
\[
  \operatorname{Ext}_{\Osh_X}^{p}(X;\Osh_Y,\Omega_X^{p})
  \longrightarrow
  \operatorname{Ext}_{\Osh_X}^{p}(X;\Osh_X,\Omega_X^{p})
  =H^{p}(X,\Omega_X^{p}).
\]
그리하여 $H^{p}(X,\Omega_X^{p})$의 한 원소를 얻으며, 이를 $P_X(Y)$로 나타내고
\emph{$Y$가 $X$ 안에서 갖는 코호몰로지류}라 부른다. 따라서 이 원소는 단면
$s_Y$가 $\Extsh_{\Osh_X}^{p}(\Osh_Y,\Omega_X^{p})$에서 정하는 원소로부터 다음
준동형 도식을 통해 얻어진다:
\begin{equation}\tag{4.3}
\begin{aligned}
H^{p}(X,\Omega_X^{p})
&=\operatorname{Ext}_{\Osh_X}^{p}(X;\Osh_X,\Omega_X^{p})\\
&\longleftarrow\operatorname{Ext}_{\Osh_X}^{p}(X;\Osh_Y,\Omega_X^{p})
\xrightarrow{\ \sim\ }
H^{0}\!\left(X,\Extsh_{\Osh_X}^{p}(\Osh_Y,\Omega_X^{p})\right).
\end{aligned}
\end{equation}

차원 $n-p$인 비특이 사이클이라 함은 차원 $n-p$인 $X$ 안의 비특이 기약
부분다양체들로 생성되는 자유가환군의 한 원소를 뜻한다. 그러면 함수
$Y\longmapsto P(Y)$는 차원 $n-p$인 $X$ 위의 비특이 사이클들의 군에서
군 $H^{p}(X,\Omega_X^{p})$으로 가는 준동형으로 연장된다.

$Z^{n-p}$와 $Z'^{n-p'}$를 각각 차원이 $n-p$와 $n-p'$인 비특이 사이클이라
하자. $Z$의 모든 성분이 $Z'$의 모든 성분과 횡단적으로 교차하면, 두 사이클이
\emph{횡단적으로 교차한다}고 한다. 그러면 사이클 $Z.Z'$가 정의되며, 이는
차원이 $n-p-p'$인 비특이 사이클이다. 이때 다음 정리가 성립한다.

\src{179}
\result{정리}{1} $Z^{n-p}$와 $Z'^{n-p'}$가 횡단적으로 교차하는 비특이
사이클이면 다음이 성립한다.
\begin{equation}\tag{4.4}
  P_X(Z.Z')=P_X(Z).P_X(Z'),
\end{equation}
여기서 오른쪽 항의 곱은 컵곱
\[
  H^{p}(X,\Omega_X^{p})\times H^{p'}(X,\Omega_X^{p'})
  \longrightarrow H^{p+p'}(X,\Omega_X^{p+p'})
\]
이다. (여기서는 $X$가 어떤 사영 공간의 국소 닫힌 부분집합과 동형이라고
가정한다.)

마지막 가정은 $X$ 위의 모든 연접 대수적 층이 국소 자유인 연접 대수적
층의 몫임을 결론내리는 데에만 사용된다(세르). 따라서 그러한 층은
국소 자유층에 의한 왼쪽 분해를 갖는다. 정리 1을 증명하기 위해
$Z$와 $Z'$가 횡단적으로 교차하는 기약 비특이 부분다양체 $Y$와 $Y'$라고
가정해도 된다. $L_{*}$를 $\Osh_Y$의 국소 자유층에 의한 왼쪽
분해라 하자. 그러면 명제 3에 따라 준동형 그림 (4.3)은 다음 그림과
동일시된다.
\[
R^{p}\Gamma(\Omega_X^{p})
\xleftarrow{\alpha}
\mathbb R^{p}\Gamma\!\left(\Homsh_{\Osh_X}(L_{*},\Omega_X^{p})\right)
\xrightarrow[\sim]{\beta}
\Gamma\!\left(H^{p}\bigl(\Homsh_{\Osh_X}(L_{*},\Omega_X^{p})\bigr)\right).
\]
여기서 $\Gamma$는 $X$ 위의 아벨군의 층들의 범주에서 ``단면군''을 취하는
함자이고, $\mathbb R^{p}\Gamma$는 이 함자의 $p$차 초코호몰로지를,
$R^{p}\Gamma$는 그 $p$번째 유도함자를 나타낸다. 단순히 하기 위해
$L_0=\Osh_X$이고 증대 준동형 $L_0\longrightarrow\Osh_Y$가 자연스러운
준동형이라고 가정해도 된다. 그러면 $\alpha$는 복합체의 준동형
$\Osh_X\longrightarrow L_{*}$로부터 얻어진다. 여기서 $\Osh_X$는 차수
$0$에만 놓인 복합체로 간주하며, 층들의 복합체 $K$가 차수 $0$에만
놓이면 $\mathbb R^{p}\Gamma(K)=R^{p}\Gamma(K_0)$임을 고려한다. 준동형
$\beta$는 잘 알려진 ``가장자리 사상(edge-homomorphism)''이다.
$\Osh_{Y'}$의 국소 자유층에 의한 분해 $L'_{*}$에 관한 유사한 그림을 생각하고,
다음과 같이 짝짓기로 이루어진 가환 그림을 생각하자.
\begin{equation}\tag{4.5}
\resizebox{\textwidth}{!}{$
\begin{array}{ccc}
R^{p}\Gamma(\Omega_X^{p})
&\longleftarrow
\mathbb R^{p}\Gamma\!\left(\Homsh_{\Osh_X}(L_{*},\Omega_X^{p})\right)
&\xrightarrow{\sim}
\Gamma H^{p}\!\left(\Homsh_{\Osh_X}(L_{*},\Omega_X^{p})\right)
\\[3pt]
\times&\times&\times\\[3pt]
R^{p'}\Gamma(\Omega_X^{p'})
&\longleftarrow
\mathbb R^{p'}\Gamma\!\left(\Homsh_{\Osh_X}(L'_{*},\Omega_X^{p'})\right)
&\xrightarrow{\sim}
\Gamma H^{p'}\!\left(\Homsh_{\Osh_X}(L'_{*},\Omega_X^{p'})\right)
\\[3pt]
\downarrow&\downarrow&\downarrow\\[3pt]
R^{p+p'}\Gamma(\Omega_X^{p+p'})
&\longleftarrow
\mathbb R^{p+p'}\Gamma\!\left(
\Homsh_{\Osh_X}(L_{*}\otimes L'_{*},\Omega_X^{p+p'})\right)
&\xrightarrow{\sim}
\Gamma H^{p+p'}\!\left(
\Homsh_{\Osh_X}(L_{*}\otimes L'_{*},\Omega_X^{p+p'})\right).
\end{array}
$}
\end{equation}

\src{180}
오른쪽 두 열의 짝짓기는 층 복합체들의 짝짓기
\[
\Homsh_{\Osh_X}(L_{*},\Omega_X^{p})
\times\Homsh_{\Osh_X}(L'_{*},\Omega_X^{p'})
\longrightarrow
\Homsh_{\Osh_X}(L_{*}\otimes L'_{*},\Omega_X^{p+p'}),
\]
으로부터 얻어진다. 이 짝짓기는 외적
$\Omega_X^{p}\times\Omega_X^{p'}\longrightarrow\Omega_X^{p+p'}$을 사용하여
정의한다. 첫째 열의 짝짓기는 (외적에 관한) 컵곱이다. 나는 (4.5)의
마지막 행이, (4.3)과 유사하되 $Y$를 $Y\cap Y'$으로, $p$를 $p+p'$으로
바꾼 동형들의 그림과 동일시된다고 주장한다. 이를 위해서는
$L\otimes L'$이 $\Osh_{Y\cap Y'}$의 (명백히 국소 자유인) 분해임을
보이면 충분하다. 실제로
\[
  H_0(L\otimes L')=\Osh_Y\otimes\Osh_{Y'}=\Osh_{Y\cap Y'}
\]
이고
\[
  H_i(L\otimes L')=\operatorname{Tor}_{i}^{\Osh_X}(\Osh_Y,\Osh_{Y'})=0
  \qquad(i>0),
\]
인데, 이는 $Y$와 $Y'$가 횡단적으로 교차하기 때문이다. 이제 다음 식에서
정리 1이 따른다.
\begin{equation}\tag{4.6}
  s_Y.s_{Y'}=s_{Y.Y'}
\end{equation}
(여기서 왼쪽 항의 곱은 (4.5)의 마지막 열에 나타난 곱이다.) 순전히
국소적인 식 (4.6)은 $L_{*}$와 $L'_{*}$로 명제 4에서 고려한 분해들을
택하면 어렵지 않게 확인된다. 마찬가지로 (더 쉬운 증명으로)
$Z\longmapsto P_X(Z)$가 데카르트 곱과 양립함을 보인다.
\begin{equation}\tag{4.7}
  P_{X\times X'}(Z\times Z')=P_X(Z)\otimes P_{X'}(Z')
\end{equation}
(이 식은 $Z$와 $Z'$가 각각 비특이 다양체 $X$와 $X'$ 위의 비특이
사이클이고, $Z\times Z'$를 $X\times X'$ 위의 비특이 사이클로 간주할 때
성립한다.) (4.4)와 (4.7)에서 $P_X(Z)$가 비특이 다양체의 사상
$f:X\longrightarrow X'$에 의한 ``역상'' 연산과도 양립함이 따른다.
\begin{equation}\tag{4.8}
  P_X\bigl(f^{-1}(Z')\bigr)=f^{*}\bigl(P_{X'}(Z')\bigr),
\end{equation}
이 식은 $Z'$가 $X'$ 위의 비특이 사이클이고 $f$가 $Z'$에 ``횡단적''일 때,
즉 $f$의 그래프가 $X\times X'$ 안의 사이클 $X\times Z'$와 횡단적으로
교차할 때 성립한다.

\src{181}
\result{따름정리}{1} $X$, $X'$을 사영공간 안에서 국소 닫힌 두 비특이
다양체라 하고, $X'$은 완비라고 하자. $U$를 $X\times X'$ 위의 비특이
사이클이라 하고, $a$, $b$를 $X'$의 두 점이라 하되 $U$가 사이클
$X\times(a)$ 및 $X\times(b)$와 횡단적으로 교차한다고 하자. $Z$, $Z'$을
다음을 만족하는 $X$ 위의 비특이 사이클이라 하자.
\[
  Z\times(a)=(X\times(a)).U,
  \qquad
  Z'\times(b)=(X\times(b)).U.
\]
그러면
\[
  P_X(Z)=P_X(Z').
\]

실제로 $f_a:X\longrightarrow X\times X'$를 $f_a(x)=(x,a)$로 정의하자.
그러면 (4.8)에 따라 $P(Z)=f_a^{*}(P(U))$이고, 마찬가지로
$P(Z')=f_b^{*}(P(U))$이다. 한편 K\"unneth 공식
\[
  H^{*}(X\times X',\Omega_{X\times X'}^{*})
  =H^{*}(X,\Omega_X^{*})\otimes H^{*}(X',\Omega_{X'}^{*})
\]
과 $H^{0}(X',\Osh_{X'})$가 스칼라들로만 이루어져 있다는 사실을
이용하면 $f_a^{*}=f_b^{*}$를 쉽게 얻으며, 이로써 결론이 따른다.

각 $x\in X$에 대하여 점 $(x)$는 $X$의 여차원 $n$인 비특이
부분다양체이므로 $H^{n}(X,\Omega_X^{n})$의 원소 $\varepsilon_x$를 정한다.
$X$가 비특이 사영 다양체이면 따름정리 1에 따라 $\varepsilon_x$는 선택한
점 $x$에 의존하지 않는다. 이를 $\varepsilon_X$로 나타내고
$H^{n}(X,\Omega_X^{n})$의 \emph{기본류}라 부른다.

\result{비고}{} 만족스러운 이론을 얻으려면 임의의 사이클 $Z$에 대하여
$P_X(Z)$를 정의하고, 사이클들의 고유 교차에 대하여 정리 1을 증명해야 할
것이다. (이 강연록을 쓸 당시에는 아직 일반적인 경우에 이 일이 이루어지지
않았다.) 이것이 이루어졌다고 가정하면 따름정리 1은 다음과 같이 된다.
$Z$와 $Z'$이 대수적으로 동치인 두 사이클이면 $P_X(Z)=P_X(Z')$이다
(이 명제는 $Z$와 $Z'$이 비특이라 하더라도 앞의 결과에서 나오는 것 같지
않다).

\fsection{5}{쌍대성 정리}

이 절에서 $X$는 차원 $n$인 비특이 사영 다양체를 나타낸다.

\result{정리}{2} $H^{n}(X,\Omega_X^{n})$의 기본류 $\varepsilon_X$는 이
벡터공간의 기저를 이룬다.

(증명은 아래에서 제시한다.) 이 정리에 따라
$H^{n}(X,\Omega_X^{n})$를 체 $k$와 동일시할 수 있다. 이제 제2절에서
고찰한 짝짓기들을 생각하자. 이들은 특히 다음 짝짓기를 준다.

\src{182}
\[
\operatorname{Ext}_{\Osh_X}^{p}(X;\Osh_X,F)
\times\operatorname{Ext}_{\Osh_X}^{n-p}(X;F,\Omega_X^{n})
\longrightarrow
\operatorname{Ext}_{\Osh_X}^{n}(X;\Osh_X,\Omega_X^{n}),
\]
즉,
\begin{equation}\tag{5.1}
  H^{p}(X,F)
  \times\operatorname{Ext}_{\Osh_X}^{n-p}(X;F,\Omega_X^{n})
  \longrightarrow H^{n}(X,\Omega_X^{n}).
\end{equation}
정리 2를 고려하면 이 짝짓기는 다음 준동형을 정한다.
\begin{equation}\tag{5.2}
  \operatorname{Ext}_{\Osh_X}^{n-p}(X;F,\Omega_X^{n})
  \longrightarrow\bigl(H^{p}(X,F)\bigr)'.
\end{equation}
이 준동형은 $F$에 관해 함자적이며, 완전열
$0\longrightarrow F'\longrightarrow F\longrightarrow F''
\longrightarrow0$에 딸린 연결 준동형들과 가환한다.

\result{정리}{3} 준동형 (5.2)는 동형이다.

특히 세르의 다음 결과를 다시 얻는다.

\result{따름정리}{} $E$를 $X$ 위의 대수적 벡터 다발이라 하고,
$\Osh_X(E)$를 $E$의 정칙 단면의 싹들로 이루어진 층이라 하자. 그러면
다음 표준 동형들이 성립한다.
\begin{equation}\tag{5.3}
  \bigl(H^{p}(X,\Osh_X(E))\bigr)'
  =H^{n-p}\bigl(X,\Omega_X^{n}\otimes\Osh_X(E')\bigr).
\end{equation}
명제 1에 정리 3과 따름정리 1을 적용하면 충분하다.

정리 2와 정리 3은 다음 진술로부터 따른다.

\noindent\textbf{(D) 다음 준동형}
\begin{equation}\tag{5.2 bis}
  \operatorname{Ext}_{\Osh_X}^{n-p}(X;F,\Omega_X^{n})
  \longrightarrow\bigl(H^{p}(X,F)\bigr)'\otimes L,
  \qquad L=H^{n}(X,\Omega_X^{n}),
\end{equation}
은 짝짓기 (5.1)에서 유도되며 동형이다.

실제로 (D)가 정리 2를 함의함을 보이자. $k_x=\Osh_{(x)}$를 한 점
$x\in X$로 이루어진 축소 다양체의 국소환의 층이라 하자. 표준 준동형
$\Osh_X\longrightarrow k_x$와 이에 딸린 준동형
\begin{equation}\tag{5.3 bis}
  H^{0}(X,\Osh_X)\longrightarrow H^{0}(X,k_x),
\end{equation}
을 생각하자. 그 쌍대 준동형은 다음 준동형과 동일시된다.
\begin{equation}\tag{5.4}
  \operatorname{Ext}_{\Osh_X}^{n}(X;k_x,\Omega_X^{n})\otimes L'
  \longrightarrow
  \operatorname{Ext}_{\Osh_X}^{n}(X;\Osh_X,\Omega_X^{n})\otimes L'
\end{equation}
이는 $\Osh_X\longrightarrow k_x$에 의해 $\operatorname{Ext}^{n}$에
유도되는 준동형
\begin{equation}\tag{5.5}
  \operatorname{Ext}_{\Osh_X}^{n}(X;k_x,\Omega_X^{n})
  \longrightarrow
  \operatorname{Ext}_{\Osh_X}^{n}(X;\Osh_X,\Omega_X^{n})
  =H^{n}(X,\Omega_X^{n}).
\end{equation}
에서 유도된다.

\src{183}
(5.3 bis)가 동형이므로\footnote{인쇄본은 182쪽에서 두 번째 공식
“(5.3)”을 제시하고 여기서 그것을 참조한다. 식별자 “(5.3 bis)”는 이
공식을 첫 번째 (5.3)과 구별한다.} (5.4)도 동형이고, 따라서 (5.5)도
동형이다. 한편 (4.2)에 따라 $s_{(x)}$는
$\operatorname{Ext}_{\Osh_X}^{n}(X;k_x,\Omega_X^{n})$의 기저이므로, 그
상인 $\varepsilon_X$는 $H^{n}(X,\Omega_X^{n})$의 기저이다.

이제 진술 (D)를 증명하는 일만 남았다. 이는 다음 보조정리들에 요약된
초등적 사실들로부터 순전히 형식적으로 따를 것이다. 이 보조정리들에서는
$X$를 차원 $r$인 사영공간 $P$의 닫힌 부분집합으로 가정한다. $X$는
특이일 수도 있고 비특이일 수도 있다. [3]에서 $\Osh(m)$으로 나타낸 $P$
위의 층을 $\Osh_P(m)$으로 나타내고, $X$ 위의 그와 같은 층을
$\Osh_X(m)$으로 나타낸다.

\result{보조정리}{2} $X=P$이고 $F=\Osh_P(m)$이면 진술 (D)가 성립한다.

이 보조정리는 직접 계산으로 확인된다. $H^{i}(P,\Osh_P(m))$의 명시적
계산은 [3]에 있지만 더 초등적으로도 할 수 있다. 짝짓기 (5.1)를
명시하는 데 필요한 컵곱
\[
  H^{i}(P,\Osh_P(m))\times H^{j}(P,\Osh_P(m'))
  \longrightarrow H^{i+j}(P,\Osh_P(m+m'))
\]
의 계산에는 어려움이 없다.\footnote{인쇄본에는 $\Osh_P(r+r')$로 되어
있으나, 두 인자 $\Osh_P(m)$과 $\Osh_P(m')$에 따라
$\Osh_P(m+m')$이어야 한다.}

\result{보조정리}{3} $X$ 위의 모든 연접 대수적 층 $F$는
$\Osh_X(-m)^{k}$의 어떤 몫층과 동형이다. 여기서 $m$은 원하는 만큼
크게 잡을 수 있다.

$m$이 충분히 클 때 $F\otimes\Osh_X(m)$이 ``그 단면들로 생성된다''는
사실에서 따른다([3] 참조).

\result{보조정리}{4} $i>0$이라 하자. 그러면
$m$이 충분히 클 때 $H^{r-i}(P,\Osh_P(-m))=0$이고, $X$ 위의 임의의
연접 대수적 층 $B$에 대하여 다음이
\[
  \operatorname{Ext}_{\Osh_X}^{i}(X;\Osh_X(-m),B)=0
\]
$m$이 충분히 클 때 성립한다.

첫 번째 명제는 위에서 언급한 명시적 계산으로부터 따른다. 두 번째 명제에
대해서는 다음 동형이 있음을 주목한다.
\[
  \operatorname{Ext}_{\Osh_X}^{i}(X;\Osh_X(-m),B)
  =H^{i}(X,B\otimes\Osh_X(m))
\]
(명제 1, 따름정리 1). 따라서 [3]의 잘 알려진 결과에 의해 결론이
따른다. 앞의 두 보조정리를 결합하면 다음 따름정리를 얻는다.

\result{따름정리}{} $i>0$이라 하자. 그러면 $P$ 위의 연접 대수적 층의
범주에서 정의된 함자 $F\longmapsto H^{r-i}(P,F)$는 에피사상에 의해
소거 가능하고, $X$ 위의 연접 대수적 층의 범주에서 정의된 함자
$\operatorname{Ext}_{\Osh_X}^{i}(X;F,B)$도 마찬가지이다.

\src{184}
\result{보조정리}{5} $A$와 $B$를 $X$ 위의 두 연접 대수적 층이라 하고,
$A(m)=A\otimes\Osh_X(m)$이라 놓자. 그러면 $m$이 충분히 클 때 표준
준동형
\[
  \operatorname{Ext}_{\Osh_X}^{1}(X;A(-m),B)
  \longrightarrow H^{0}\bigl(X,\Extsh_{\Osh_X}^{1}(A(-m),B)\bigr)
  =H^{0}\bigl(X,\Extsh_{\Osh_X}^{1}(A,B)(m)\bigr)
\]
은 동형이다.

이는 명제 1의 스펙트럴 수열을 $A(-m)$과 $B$에 적용하면 곧바로 따른다.
실제로 다음을 얻기 때문이다.
\[
  E_{2}^{p,q}(A(-m),B)
  =H^{p}\bigl(X,\Extsh_{\Osh_X}^{q}(A(-m),B)\bigr)
  =H^{p}\bigl(X,\Extsh_{\Osh_X}^{q}(A,B)(m)\bigr),
\]
$p>0$이고 $m$이 충분히 크면 이는 0이다.

이제 $X=P$인 경우에 (D)를 증명하자. 먼저 $p=n$일 때 (5.2 bis)가
동형임을 보인다. 이때 양변은 모두 왼쪽 완전 함자이다
($H^{r+1}(P,F)=0$이기 때문이다). 따라서 보조정리 3에 의해
$F=\Osh_P(-m)$인 경우에 명제를 증명하면 충분한데, 이 경우는
보조정리 2에 포함되어 있다. 준동형들 (5.2 bis)는 함자적이고
연결 준동형들과 양립하며, $p<n$일 때 (5.2 bis)의 양변은 모두
$F$에 관하여 에피사상에 의해 소거 가능한 함자이므로(보조정리 4의 따름정리),
표준적인 논법에 따라 (5.2 bis)는 모든 $p$에 대하여 동형이다.
이로써 사영공간에 대한 쌍대성 정리가 증명되었다. 이제 $X$는 임의이되
비특이라고 하자. $P$에 대한 쌍대성 정리에 의해 다음 동형이 있다.
\[
  H^{n}(X,F)'=H^{n}(P,F)'
  =\operatorname{Ext}_{\Osh_P}^{r-n}(P;F,\Omega_P^{r}).
\]
4절의 보조정리 1에 의해 마지막 항은 다음과 동일시된다.
\[
  \operatorname{Hom}_{\Osh_P}(P;F,\omega_X^{n})
  =\operatorname{Hom}_{\Osh_X}(X;F,\Omega_X^{n})
  =\operatorname{Ext}_{\Osh_X}^{0}(X;F,\Omega_X^{n}),
\]
따라서 다음 동형을 얻는다.
\begin{equation}\tag{5.6}
  H^{n}(X,F)'=\operatorname{Hom}_{\Osh_X}(X;F,\Omega_X^{n})
  =\operatorname{Ext}_{\Osh_X}^{0}(X;F,\Omega_X^{n}).
\end{equation}
$F=\Omega_X^{n}$으로 놓으면 다음 동형을 얻는다.
\begin{equation}\tag{5.7}
  \gamma:H^{n}(X,\Omega_X^{n})\xrightarrow{\sim}k.
\end{equation}

\src{185}
동형 (5.6)은 (5.7)을 사용하여 $L=k$로 놓았을 때 $p=n$에 대한
(5.2 bis)와 다름없음을 확인할 수 있다. 따라서 (5.2 bis)는 $p=n$일 때
동형이다. 모든 $p$에 대하여 동형임을 증명하려면, $p<n$일 때
(5.2 bis)의 양변이 $F$에 관하여 에피사상에 의해 소거 가능한 함자임을
보이면 충분하고, 보조정리 3을 고려하면 더 강하게 양변이
$F=\Osh_X(-m)$이고 $m$이 충분히 클 때 0임을 보이면 충분하다. 첫 번째
변에 대해서는 보조정리 4로부터 따르고, 두 번째 변에 대해서는
$P$ 위의 쌍대성 정리를 사용하여 다음과 같이 쓴다.
\[
  H^{p}(X,\Osh_X(-m))'
  =\operatorname{Ext}_{\Osh_P}^{r-p}
    (P;\Osh_X(-m),\Omega_P^{r}).
\]
두 번째 변은 $p<n$이고 $m$이 충분히 클 때 0이다. 이는 보조정리 5에서
$X=P$로 놓고, $P$ 위의 연접 대수적 층으로서 $\Osh_X$가 코호몰로지
차원 $\leq r-n$을 갖는다는 사실($X$가 비특이이므로)을 사용하면 따른다.
따라서
\[
  \Extsh_{\Osh_P}^{r-p}(\Osh_X,\Omega_P^{r})=0
  \qquad\text{$p<n$일 때}.
\]
\fsection{6}{특이 다양체에 대한 쌍대성 정리}

$X$를 차원이 $r$인 사영공간 $P$의 차원이 $n$인 닫힌 부분집합이라
하자. 여기서는 공식 (5.6)을 다음과 같이 써야 한다.
\begin{equation}\tag{6.1}
  H^{n}(X,F)'
  \simeq\operatorname{Hom}_{\Osh_X}(X;F,\omega_X^{n})
  =\operatorname{Ext}_{\Osh_X}^{0}(X;F,\omega_X^{n}),
\end{equation}
여기서 다음과 같이 놓았다.
\begin{equation}\tag{6.2}
  \omega_X^{n}=E^{0}(\Osh_X)
  =\Extsh_{\Osh_P}^{r-n}(\Osh_X,\Omega_P^{r}).
\end{equation}
명제 5에서 말했듯이, 이렇게 도입한 층은 실제로 $X$를 비특이 다양체
$P$에 넣는 데 선택한 몰입 사상에 의존하지 않는다.
(6.1)에서 $F=\omega_X^{n}$로 놓으면 다음을 얻는다.
\begin{equation}\tag{6.2 bis}
  H^{n}(X,\omega_X^{n})'
  \simeq\operatorname{Hom}_{\Osh_X}(X;\omega_X^{n},\omega_X^{n}),
\end{equation}
\footnote{인쇄본도 이 공식에 ``(6.2)''라는 번호를 붙였다. 접미사
``bis''는 두 출현을 구별한다.} 따라서 다음에 대응하는
$H^{n}(X,\omega_X^{n})'$의 특별한 원소가 존재한다.
$\omega_X^{n}$에서 자기 자신으로 가는 항등 준동형이다:
\begin{equation}\tag{6.3}
  \gamma:H^{n}(X,\omega_X^{n})\longrightarrow k.
\end{equation}

\src{186}
이제 $\operatorname{Ext}$의 합성으로 정의되는 짝지음들을 생각하자:
\begin{equation}\tag{6.4}
  H^{p}(X,F)\times
  \operatorname{Ext}_{\Osh_X}^{n-p}(X;F,\omega_X^{n})
  \longrightarrow H^{n}(X,\omega_X^{n}),
\end{equation}
이를 (6.3)의 준동형 $\gamma$와 합성하면 경계 준동형과
양립하는 다음 함자적 준동형들을 얻는다:
\begin{equation}\tag{6.5}
  \operatorname{Ext}_{\Osh_X}^{n-p}(X;F,\omega_X^{n})
  \longrightarrow H^{p}(X,F)'
\end{equation}
(이는 (5.2)를 일반화한다). $p=n$일 때에는 이렇게 하여 (6.1)의
동형사상을 얻음을 확인할 수 있다. 이제 다음 정리가 성립한다.

\result{정리}{3 bis} 주어진 정수 $k\geq0$에 대하여 $X$에 관한
다음 네 조건은 서로 동치이다:
\begin{enumerate}
  \item[\textit{i.}] 함자적 준동형 (6.5)는 $n-k\leq p\leq n$일 때
    동형사상이다.
  \item[\textit{ii.}] 충분히 큰 $m$과 $n-k\leq p<n$에 대하여
    $H^{p}(X,\Osh_X(-m))=0$이다.
  \item[\textit{iii.}] $X$ 위의 연접 대수적 층의 범주에서 함자
    $H^{p}(X,F)$는 $n-k\leq p<n$일 때 여소거 가능하다.
  \item[\textit{iv.}] $0<i\leq k$에 대하여
    $E^{i}(\Osh_X)=\Extsh_{\Osh_P}^{r-n+i}(\Osh_X,\Omega_P^{r})=0$
    이다.
\end{enumerate}

\noindent\textbf{증명.} --- \textit{i.}$\Rightarrow$\textit{ii.}는
보조정리 4에서, \textit{ii.}$\Rightarrow$\textit{iii.}은 보조정리 3에서,
\textit{iii.}$\Rightarrow$\textit{i.}은 잘 알려진 표준 논증에서 따른다.
실제로 그때 (6.5)의 양변은 $n-k\leq p<n$에서 여소거 가능한 함자이고
(첫째는 보조정리 4에 따라 그러하기 때문이다). 끝으로
\textit{ii.}$\Leftrightarrow$\textit{iv.}를 증명하자. 이는 다음 명제의
따름정리에서 나온다:

\result{명제}{6} $F$를 $X$ 위의 연접 대수적 층이라 하고,
$i$를 정수라 하자. 그러면 충분히 큰 $m$에 대하여 다음 동형사상이 성립한다:
\begin{equation}\tag{6.6}
  H^{i}(X,F(-m))'
  \simeq H^{0}\bigl(X,E^{n-i}(F)(m)\bigr),
\end{equation}
여기서 다음과 같이 놓는다(n\textsuperscript{o} 3, 명제 5와 비교):
\begin{equation}\tag{6.7}
  E^{j}(F)=\Extsh_{\Osh_P}^{r-n+j}(F,\Omega_P^{r}).
\end{equation}
실제로 $P$에 대한 쌍대성 정리에 의해 (6.6)의 좌변은
$\operatorname{Ext}_{\Osh_P}^{r-i}(P;F(-m),\Omega_P^{r})$와 동형이므로,
(6.6)은 보조정리 5에서
따른다.

\src{187}
\result{따름정리}{} 충분히 큰 $m$에 대하여 $H^{i}(X,F(-m))=0$이
성립할 필요충분조건은 $E^{n-i}(F)=0$인 것이다.

$E^{j}(F)$들은 실제로 고려 중인 사영 몰입 사상에 의존하지 않음을 상기하자.
따름정리의 조건은 순전히 국소적이다. 따라서
이 조건이 $F$에 대해 성립하면 국소적으로
어떤 층 $F^{n}$과 동형인 모든 층에 대해서도 성립한다. 특히 이 조건이
$\Osh_X$에 대해 성립하면 모든 국소 자유 연접 대수적 층에 대해
성립한다. 예를 들어 $X$가 비특이면 모든 $i<n$에 대해 그러하고,
$X$의 어느 성분도 한 점만으로 이루어지지 않으면 $i=0$에 대해,
$X$가 정규이고 모든 성분의 차원이 $>1$이면 $i=0,1$에 대해 그러하다(참조
[3]).\footnote{인쇄본은 여기서 ``$S$가 정규이다''라고 되어 있지만, 이 단락은
다양체 $X$만을 다룬다.} 모든 $i<k$에 대해 그러할 필요충분조건은
정의에 따라 국소환
$\Osh_x$ ($x\in X$)들이 ``호몰로지적 여차원'' $\geq k$를 갖는 것이다
(이 개념은 [4] 참조). $k=n$이면 정리 3 bis에 의해 이는
쌍대성 정리가 $X$에 대해 성립한다는 것, 즉 (6.5)가
모든 $p$와 모든 $F$에 대해 동형사상이라는 뜻이다. 이에 대한 국소환
$\Osh_x$의 동치 조건을 많이 제시할 수 있는데(Nagata), 예를 들어
코언--매콜리의 등차원성 정리를 만족하는 조건이 그러하다. 이는
예를 들어 $X$가 국소적으로 비특이 주위 다양체 안의 ``완전교차''이면
성립한다.

\fsection{7}{푸앵카레 쌍대성}

$X$를 차원 $n$인 비특이 사영 다양체라 하자. 그러면
\[
  H^{**}(X)=H^{*}(X,\Omega_X^{*})
\]
는 유한 차원의 반가환 이중등급 대수이며, 우리는 여기에 전체 차수에
따른 등급을 준다. 따라서 $H^{p,q}(X)=H^{p}(X,\Omega_X^{q})$의 차수는
$p+q$이다. $H^{*}(X)$에서 나타나는 차수들은 $0$에서 $2n$까지이다.
정리 2와 정리 3의 따름정리에 따라 $H^{*}(X)$는 차원 $2n$인
“푸앵카레 대수”이다. 즉 $H^{2n}(X)$에는 기초체 $k$와의 동형사상이
주어져 있고, 곱
\[
  H^{m}(X)\times H^{2n-m}(X)\longrightarrow H^{2n}(X)=k
\]
은 $H^{m}(X)$와 $H^{2n-m}(X)$ 사이의 쌍대성을 이룬다. 또 $Y$가
다른 비특이 사영 다양체이면, 연접 대수적 층에 대한 퀸네트 공식은
다음을 준다.
\begin{equation}\tag{7.1}
  H^{*}(X\times Y)=H^{*}(X)\otimes H^{*}(Y),
\end{equation}
을 얻으며, 이 동형사상은 푸앵카레 대수 구조들과 양립한다. 또한
$H^{*}(X)$는 가환 대수로서 $X$에 관한 반변 함자이고, 사상
$f:Y\longrightarrow X$는 명백한 방식으로 등급 대수의 준동형

\src{188}
\begin{equation}\tag{7.2}
  f^{*}:H^{*}(X)\longrightarrow H^{*}(Y).
\end{equation}
을 정의한다. 이들이 푸앵카레 대수이므로, 전치를 취하여 다음 벡터
공간의 준동형을 얻는다:
\begin{equation}\tag{7.3}
  f_{*}:H^{*}(Y)\longrightarrow H^{*}(X).
\end{equation}
4절에서, 비특이 순환에 대응하는 코호몰로지류에 대한 $f^{*}$의
작용은 그 순환들의 역상에 대응하는 코호몰로지류를 취하는 것으로
기하학적으로 해석됨을 보았다. 여기서는 (7.3)이 이와 마찬가지로
순환의 “직상” 연산에 대응함을 보이는 것이 중요하다. 적어도 적절한
비특이성 조건 아래에서는, 이는 다음 특수한 경우로부터 따른다:

\result{정리}{4} $f$가 $X^{n}$ 안의 비특이 부분다양체 $Y^{m}$에서
$X^{n}$으로 가는 포함 사상이고 $1_Y$가 $H(Y)$의 단위원을 나타내면,
다음이 성립한다.
\begin{equation}\tag{7.4}
  f_{*}(1_Y)=P_X(Y),
\end{equation}
여기서 우변은 $Y$에 대응하는 $X$의 코호몰로지류이다.

이 식은 다음과 동치이다.
\begin{equation}\tag{7.4 bis}
  \bigl\langle\xi^{m,m},P_X(Y)\bigr\rangle\varepsilon_Y
  =f^{*}(\xi^{m,m})
  \qquad
  \bigl(\xi^{m,m}\in H^{m}(X,\Omega_X^{m})\bigr),
\end{equation}
\footnote{인쇄본에는 $f_{*}(\xi^{m,m})$라고 되어 있다. 그런데
$\xi^{m,m}\in H^{m}(X,\Omega_X^{m})$이고 양변이
$H^{m}(Y,\Omega_Y^{m})$에 속하므로 필요한 사상은 $f^{*}$이다.} 여기서
$\varepsilon_Y$는 $H^{m}(Y,\Omega_Y^{m})$의 기본 원소이다.
비특이 사영 다양체의 경우 이 식은 $Y$에 대응하는 코호몰로지류의 새로운
정의를 준다. 정리 4를 증명하기 위해, 정리 3을 이용하여 준동형
\[
  H^{m}(X,\Omega_X^{m})
  \longrightarrow H^{m}(Y,\Omega_Y^{m})
  =H^{m}(X,\Omega_Y^{m})
\]
의 전치를 다음 준동형으로 해석한다.
\begin{equation}\tag{7.5}
\begin{split}
  H^{n-m}(X,\Omega_X^{n-m})
  &\simeq
  \operatorname{Ext}_{\Osh_X}^{n-m}(X;\Omega_X^{m},\Omega_X^{n})\\
  &\longleftarrow
  \operatorname{Ext}_{\Osh_X}^{n-m}(X;\Omega_Y^{m},\Omega_X^{n})
  \simeq
  \operatorname{Hom}_{\Osh_X}(X;\Omega_Y^{m},\Omega_Y^{m}).
\end{split}
\end{equation}
$H^{m}(Y,\Omega_Y^{m})$의 쌍대에 속하는 원소 $1_Y$는 우변에서
$\Omega_Y^{m}$의 항등 자기준동형에 대응하는 원소와 동일시됨을
확인할 수 있다. 또한 이 원소의 $H^{n-m}(X,\Omega_X^{n-m})$에서의
상은 바로 $P_X(Y)$임도 확인할 수 있다. 이 양립성들은 이미 4절에서
제시할 수도 있었으며, 실제로 임의의 비특이 다양체에 대하여도
서술할 수 있고 성립한다. 그중 둘째 양립성은

\src{189}
예컨대 다음 표준 자기준동형들의 도표가 가환한다는 사실에서 따른다:
\begin{equation}\tag{7.6}
\resizebox{\textwidth}{!}{$
\left\{
\begin{array}{ccccc}
\operatorname{Ext}_{\Osh_X}^{n-m}(X;\Omega_X^{m},\Omega_X^{n})
&\longleftarrow&
\operatorname{Ext}_{\Osh_X}^{n-m}(X;\Omega_Y^{m},\Omega_X^{n})
&\simeq&
\operatorname{Hom}_{\Osh_X}(X;\Omega_Y^{m},\Omega_Y^{m})\\[3pt]
\Big\downarrow&&&&\Big\downarrow\\[3pt]
\operatorname{Ext}_{\Osh_X}^{n-m}(X;\Osh_X,\Omega_X^{n-m})
&\longleftarrow&
\operatorname{Ext}_{\Osh_X}^{n-m}(X;\Osh_Y,\Omega_X^{n-m})
&\simeq&
\operatorname{Hom}_{\Osh_X}(X;\Omega_X^{m},\Omega_Y^{m})
\end{array}
\right.
$}
\end{equation}

이로써 방향이 주어진 콤팩트 다양체에서의 푸앵카레 쌍대성 형식론에
대한 정확한 대응물을 얻는다. 특히 정리 4로부터 $X\times X$의
대각선에 대응하는 코호몰로지류를 결정할 수 있다. 잘 알려진
논법으로부터, 예컨대 다음 렙셰츠 공식을 얻는다:

\result{정리}{5} $f$를, 그 모든 고정점의 중복도가 1인 비특이 사영
다양체 $X$의 자기준동형이라 하자. 그러면 이러한 고정점들의 개수는
$k$의 표수를 법으로,
$f$가 $H^{i}(X)$에 유도하는 자기준동형들의 대각합의 교대합과
합동이다.

$f$에 둔 이 제한은 4절의 비고에서 언급한 어려움 때문이다.
그러나 다음 사실에 유의하자. $f$가 모든 고정점의 중복도가 1인
자기준동형과 “호모토픽”이면 렙셰츠 공식은
여전히 성립한다.

\fsection{8}{쌍대성 정리의 일반화}

$X$를 다음 성질을 갖는 비특이 대수다양체라 하자. $X$ 위의 모든
연접인 대수적 층 $F$는 국소 자유인 어떤 연접인 대수적 층의 몫층과
동형이다(이는 $X$가 어떤 사영공간 안에서 국소 닫혀 있을 때 성립한다).
그러면 $X$ 위의 모든 연접인 대수적 층 $F$는 국소 자유층들로 이루어진
유한 분해 $L$을 갖는다. 또한 그러한 두 분해가 주어지면 언제나 세 번째
분해와, 증대 준동형들과 양립하는 그 세 번째 분해에서 앞의 두 분해로
가는 준동형들을 찾을 수 있다. 마찬가지로 $L$이 $F$의 유한 국소 자유
분해이고 준동형 $F'\longrightarrow F$가 주어졌다면, $F'$의 유한 국소
자유 분해 $L'$과 $F'\longrightarrow F$에 양립하는 준동형
$L'\longrightarrow L$이 존재한다. $F'\longrightarrow F$가 전사이면 이
준동형도 전사라고 가정할 수 있다. 따라서 두 정수 $r,s\geq0$가 주어지면,
$A_1,\ldots,A_r$; $B_1,\ldots,B_s$를 인수로 하는 두 코호몰로지적
다중함자를 정의할 수 있다.

\src{190}
이 인수들은 $X$ 위의 연접인 대수적 층들의 범주에서 변하며, 두 함자의
값은 각각 $X$ 위의 연접인 대수적 층들의 범주와
$H^{0}(X,\Osh_X)$ 위의 가군들의 범주에 속한다. 그 정의식은 다음과 같다.
\begin{equation}\tag{8.1}
\resizebox{\textwidth}{!}{$
\left\{
\begin{aligned}
\underline{T}_{r}^{s*}(A_1,\ldots,A_r;B_1,\ldots,B_s)
  &={}H^{*}\!\left(
    \Homsh_{\Osh_X}\!\left(
      L(A_1)\otimes\cdots\otimes L(A_r),
      L(B_1)\otimes\cdots\otimes L(B_s)
    \right)\right),\\[3pt]
T_{r}^{s*}(A_1,\ldots,A_r;B_1,\ldots,B_s)
  &={}\mathbf R^{*}\Gamma\!\left(
    \Homsh_{\Osh_X}\!\left(
      L(A_1)\otimes\cdots\otimes L(A_r),
      L(B_1)\otimes\cdots\otimes L(B_s)
    \right)\right).
\end{aligned}
\right.
$}
\end{equation}
이 식에서 $L(F)$는 연접인 대수적 층 $F$의 유한 국소 자유 분해를
나타내고, $\mathbf R^{*}\Gamma(K)$는 층 복합체 $K$에 관한 공간 $X$의
초코호몰로지를 나타낸다. $r$ 또는 $s$가 0이면 $L(A_i)$들의 텐서곱,
각각 $L(B_j)$들의 텐서곱을 $\Osh_X$로 바꾼다. 특히
$\underline{T}_{0}^{0}$와 $T_{0}^{0}$는 인수가 없는 등급함자이다.
$\underline{T}_{0}^{0}$는 0차에 집중되어 있고 그 차수에서 층
$\Osh_X$이며, $T_{0}^{0}$는 $H^{*}(X,\Osh_X)$와 같다. (8.1)의 우변들이
선택한 분해에 의존하지 않는다는 사실은 첫째 줄에 대해서는 어쨌든
명백하다(이 경우 문제는 국소적이다). 둘째 줄에 대해서는 앞의 일반적
고찰과 다음 층 복합체의 초코호몰로지 스펙트럼 열을 고려하면 이 사실을
얻는다.
\[
  K=\Homsh_{\Osh_X}\bigl(L(A_1)\otimes\cdots,
      L(B_1)\otimes\cdots\bigr),
\]
이 스펙트럼 열은 $K$에 관한 $X$의 초코호몰로지로 수렴하며, 그 초기항은
$H^{p}(X,H^{q}(K))$, 즉 다음과 같다.
\begin{equation}\tag{8.2}
  E_{2}^{p,q}
  =H^{p}\!\left(X,
    (\underline{T}_{r}^{s})^{(q)}
      (A_1,\ldots,A_r;B_1,\ldots,B_s)\right).
\end{equation}
그러므로 이 스펙트럼 열 자체도 선택한 분해에 의존하지 않으며, 그
수렴대상은 (8.1)의 좌변이다. 여러 인수 $A_i$, $B_j$에 관한 연결
연산자들도 쉽게 정의된다. 실제로 모든 완전열
$0\longrightarrow F'\longrightarrow F\longrightarrow F''\longrightarrow0$은
유한 국소 자유 복합체들로 이루어진 완전열로 분해할 수 있다.

함자계 $\underline{T}_{r}^{s*}$, 각각 $T_{r}^{s*}$에는 텐서 계산의
연산들과 비슷한 연산들을 정의한다. 그 정의는 정의식 (8.1)에서 곧바로
나온다. 이를테면 제2절에서 다룬 것을 일반화하는 다음 합성이 있다.
\begin{equation}\tag{8.3}
\resizebox{\textwidth}{!}{$
\begin{aligned}
&T_{r}^{s*}(A_1,\ldots,A_r;B_1,\ldots,B_s)
\times T_{r'}^{s'*}(A'_1,\ldots,A'_{r'};B'_1,\ldots,B'_{s'})\\
&\qquad\longrightarrow
T_{r+r'}^{s+s'*}
(A_1,\ldots,A_r,A'_1,\ldots,A'_{r'};
 B_1,\ldots,B_s,B'_1,\ldots,B'_{s'}),
\end{aligned}
$}
\end{equation}
이 합성은 명백한 결합법칙을 만족하며 함자적 준동형, 연결 준동형,
스펙트럼 열과 양립한다. 대칭 연산도 마찬가지로 있으며, 이를 명시하는
일은 독자에게 맡긴다. 또한 한 인수 $A_i$가 한 인수 $B_j$와 같을 때마다
축약 연산이 있다.

\src{191}
\begin{equation}\tag{8.4}
\resizebox{\textwidth}{!}{$
\begin{aligned}
&T_{r}^{s*}(A_1,\ldots,A_{i-1},C,A_{i+1},\ldots,A_r;
 B_1,\ldots,B_{j-1},C,B_{j+1},\ldots,B_s)\\
&\qquad\longrightarrow
T_{r-1}^{s-1*}(A_1,\ldots,\widehat{A_i},\ldots,A_r;
 B_1,\ldots,\widehat{B_j},\ldots,B_s).
\end{aligned}
$}
\end{equation}
더 나아가 한 인수 $A_i$가 국소 자유층이면, 한 $A_j$ ($j\neq i$)를
$A_j\otimes A_i$로 바꾸거나 한 $B_k$를 $B_k\otimes A_i'$로 바꾸는 대신
$A_i$를 없앨 수 있다. 여기서
$A_i'=\Homsh_{\Osh_X}(A_i,\Osh_X)$이다. 한 인수 $B_i$가 국소 자유층인
경우에도 이와 비슷한 계산 규칙이 있다. 특히 $\Osh_X$와 같은 인수는
언제나 없앨 수 있다. 모든 인수가 국소 자유이고 $B_i$들 가운데 많아야
하나만 예외라면, 방금 말한 규칙은 다음 함자적 동형을 준다.
\begin{equation}\tag{8.5}
  T_{r}^{s*}(A_1,\ldots,A_r;B_1,\ldots,B_s)
  =H^{*}\bigl(X,A_1'\otimes\cdots\otimes A_r'
    \otimes B_1\otimes\cdots\otimes B_s\bigr)
\end{equation}
(실제로 이는 $r=0$, $s=1$인 경우로 귀착되며 그 경우에는 명백하다.
초기항이 (8.2)인 스펙트럼 열을 직접 사용해도 된다.) 앞의 연산들에
대응하는 연산들을 $\underline{T}_{r}^{s}$에 대해서도 정의할 수 있다.
이렇게 도입한 여러 연산 사이의 관계는 텐서 계산의 대응하는 연산 사이의
관계와 같다.

$X$의 차원을 $n$이라 하자. 반복되는 인자들에 텐서 합성 (8.3)과
축약 (8.4)을 차례로 적용하면 다음 짝짓기를 얻는다.
\begin{equation}\tag{8.6}
\resizebox{\textwidth}{!}{$
  (T_{r}^{s})^{p}(A_1,\ldots;B_1,\ldots)
  \times
  (T_{s}^{r})^{n-p}(B_1,\ldots;A_1,\ldots,A_r\otimes\Omega_X^{n})
  \longrightarrow H^{n}(X,\Omega_X^{n}).
$}
\end{equation}

\result{정리}{6} $X$가 비특이 사영 대수다양체이면, 짝짓기 (8.6)은
쌍대성을 이룬다.

이는 제3정리의 따름정리로부터 순전히 형식적으로 나온다. 실제로 그
따름정리에 따르면, $K$가 국소 자유인 연접 대수적 층들의 복합체일 때
복합체 $K$에 대한 $X$의 초코호몰로지는 자연스러운 짝짓기
\begin{equation}\tag{8.7}
  \mathbf R^{p}\Gamma(K)
  \times\mathbf R^{n-p}\Gamma(K'\otimes\Omega_X^{n})
  \longrightarrow\mathbf R^{n}\Gamma(\Omega_X^{n})
  =H^{n}(X,\Omega_X^{n})
\end{equation}
를 통하여 복합체 $K'\otimes\Omega_X^{n}$에 대한 $X$의 초코호몰로지와
쌍대를 이룸을 쉽게 알 수 있다. 이는 첫 항이
$H^{p}(X,H^{q}(K))$인 스펙트럼 열\footnote{인쇄본에는
$H^{p}(H^{q}(X,K))$라고 되어 있으나, 첫 항은 스펙트럼 열 (8.2)의
것, 곧 $H^{p}(X,H^{q}(K))$이다.}과
$K'\otimes\Omega_X^{n}$에 대한 유사한 스펙트럼 열을 사용하면 알 수
있다. 앞의 결과에서 정의 (8.1)을 사용하면 정리 6이 따른다.

\src{192}
\subsection*{비고}

\noindent 1\textsuperscript{o} 정리 6 앞의 정의들에서는 $X$가
비특이일 필요가 없었다. 유한 분해만을 다루는 것이 필수적이지 않았기
때문이다. 그러나 $X$가 특이이면
$(\underline{T}_{r}^{s})^{p}(A_1,\ldots;B_1,\ldots)$이 연접층이라고
미리 단언할 수는 없다. 실제로 층들의 복합체
\[
  \Homsh_{\Osh_X}\bigl(L(A_1)\otimes\cdots,L(B_1)\otimes\cdots\bigr)
\]
에는 주어진 전체 차수마다 무한히 많은 성분이 생기기 때문이다.

\medskip
\noindent 2\textsuperscript{o} 식 (8.1)에서 $L(B_i)$ 가운데 하나를
$B_i$로 바꾸어도 됨은 쉽게 확인된다. 따라서 명제 3을 함께 고려하면
다음을 얻는다.
\begin{equation}\tag{8.8}
\left\{
\begin{aligned}
  \underline{T}_{1}^{1*}(A;B)&=\Extsh_{\Osh_X}^{*}(A,B),\\
  T_{1}^{1*}(A;B)&=\operatorname{Ext}_{\Osh_X}^{*}(X;A,B).
\end{aligned}
\right.
\end{equation}
특히 (8.6)에서 $r=s=1$, $A_1=\Osh_X$로 놓으면 정리 3을 다시 얻는다.
식 (8.8)은 또한 다음을 함의한다.
\[
  T_{0}^{1*}(B)=H^{*}(X,B),
  \qquad
  T_{1}^{0*}(A)=\operatorname{Ext}_{\Osh_X}^{*}(X;A,\Osh_X).
\]

\medskip
\noindent 3\textsuperscript{o} 식 (8.1)을 보면 일반적으로 함자
$\underline{T}_{r}^{s*}$와 $T_{r}^{s*}$는 양의 차수와 음의 차수의
성분을 모두 가진다. 비고 2를 사용하면 $X$의 차원이 $n$일 때
$\underline{T}_{r}^{s*}$의 영이 아닌 성분들은 $s>0$이면
$-(s-1)n$과 $rn$ 사이에, $s=0$이면 $0$과 $rn$ 사이에 있음을 알 수
있다. 한편 $T_{r}^{s*}$의 영이 아닌 성분들은 $s>0$이면
$-(s-1)n$과 $(r+1)n$ 사이에, $s=0$이면 $0$과 $(r+1)n$ 사이에 있다
(그리고 내가 잘못 보지 않았다면, $r>0$일 때에는 각각
$-(s-1)n$과 $rn$ 사이, $0$과 $rn$ 사이에 있다).

\section*{참고문헌}

\noindent[1] Cartier (Pierre). --- Les groupes
$\operatorname{Ext}^{s}(A,B)$, Séminaire A. Grothendieck : Algèbre homologique,
제1권, 1957, 제3호.

\noindent[2] Grothendieck (Alexandre). --- Sur quelques points d'algèbre
homologique, \emph{Tohoku Math. J.}, 제9권, 1957, 119--183쪽.

\noindent[3] Serre (Jean-Pierre). --- Faisceaux algébriques cohérents,
\emph{Annals of Math.}, 제61권, 1955, 197--278쪽.

\noindent[4] Serre (Jean-Pierre). --- Sur la dimension homologique des anneaux
et des modules noethériens, \emph{Proc. Intern. Symp. on Alg. Number Theory}
[1955, Tokyo et Nikko]. --- Tokyo, Science Council of Japan, 1956;
175--189쪽.

\clearpage
\src{193}
\section*{추기}

13쪽의 비고에서 지적한 어려움들은 쌍대성 정리를 임의의 대수다양체
(임의의 특이점을 가진 대수다양체)로 확장함으로써 이제 완전히
해결되었다. ``스킴 이론''의 일반적 틀에서 정식화할 수 있는 이 이론에
관해서는 현재 준비 중인 J.~Dieudonné와 A.~Grothendieck의
``Éléments de Géométrie algébrique''를 참조한다. 몇 가지 실마리는 다음
글에서 찾을 수 있다.

\begin{quote}
Grothendieck (A.). --- Cohomology theory of algebraic varieties,
\emph{Congrès int. des Mathématiciens}, 1958, Edinburgh (출간 예정).
\end{quote}

\begin{flushright}
[1959년 4월]
\end{flushright}

\begin{center}\rule{36mm}{0.5pt}\end{center}

% FGA 정본 인쇄문 보존 장치: 원자적 통합 R1.
\par\smallskip
\begingroup\footnotesize
\noindent\textit{FGA-CANON-ERR-0018. 인쇄본의 독법은 C078에 계속 기록되어 있으며, FGA-CANON-DISP-0021이 현행 본문의 최소 수정을 규율한다.}\par
\noindent\textit{FGA-CANON-ERR-0019. 인쇄본의 독법은 C079에 계속 기록되어 있으며, FGA-CANON-DISP-0022가 현행 본문의 최소 수정을 규율한다.}\par
\noindent\textit{FGA-CANON-ERR-0026. 인쇄본의 독법은 C087에 계속 기록되어 있으며, FGA-CANON-DISP-0028이 현행 본문의 최소 수정을 규율한다.}\par
\noindent\textit{FGA-CANON-ERR-0029. 인쇄본의 독법은 C092에 계속 기록되어 있으며, FGA-CANON-DISP-0031이 현행 본문의 최소 수정을 규율한다.}\par
\noindent\textit{FGA-CANON-ERR-0044. 하나의 유한 표시로부터 모든 차수에서 동형이라고 한 인쇄본의 주장은 C110에 계속 기록되어 있으며, FGA-CANON-DISP-0044가 필요한 부분 분해와 유사연접성 가설을 갖춘 현행 명제를 규율한다.}\par
\endgroup
